Man bestimme die quadratfreie Faktorisierung der folgenden Polynome
\begin{teilaufgaben}
\item $x^3-3x^2+3x-1$
\item $x^4+2x^2+1$
\item $x^5-x^4-5x^3+x^2+8x+4$
\end{teilaufgaben}

\begin{loesung}
\begin{teilaufgaben}
\item Das Polynom $q(x)=x^3-3x^2+3x-1$ hat die Ableitung $q'(x)=3x^2-6x+3$,
der grösste gemeinsame Teiler ist
\[
g_1(x)
=
\operatorname{ggT}(q(x),q'(x))
=
x^2-2x+1
\]
und weiter
\begin{align*}
g_0(x)
&=
x^3-3x^2+3x-1
&&
&&
\\
g_1(x)
&=
x^2-2x+1
&
p_1(x)
&=
x-1,
&
q_1(x)
&=
1,
\\
g_2(x)
&=
x-1,
&
p_2(x)
&=
x-1,
&
q_2(x)
&=
1,
\\
g_3(x)
&=
1,
&
p_3(x)
&=
x-1,
&
q_3(x)
&=
x-1,
\\
g_4(x)
&=
1,
&
p_4(x)
&=
1.
&&
\end{align*}
Somit ist $x^3-3x^2+3x-1=(x-1)^3$ die quadratfreie Faktorisierung.
\item
Für das Polynom $q(x)=x^4+2x^2+1$ ist 
\begin{align*}
g_0(x)
&=
x^4+2x^2+1,
&&
&&
\\
g_1(x)
&=
x^2+1,
&
p_1(x)
&=
x^2+1,
&
q_1(x)
&=
1,
\\
g_2(x)
&=
1,
&
p_2(x)
&=
x^2+1,
&
q_2(x)
&=
x^2+1,
\\
g_3(x)
&=
1,
&
p_3(x)
&=
1.
&&
\end{align*}
Die quadratfreie Faktorisierung ist daher $x^4+x^2+1 = (x^2+1)^2$.
\item
Für $q(x)=x^5-4x^4-5x^3+x^2+8x+4$ ist
\begin{align*}
g_1(x)
&=
x^3-3x-2,
&
p_1(x)
&=
x^2-x-2,
&
q_1(x)
&=
1,
\\
g_2(x)
&=
x+1,
&
p_2(x)
&=
x^2-x-2,
&
q_2(x)
&=
x-2,
\\
g_3(x)
&=
1,
&
p_3(x)
&=
x+1,
&
q_3(x)
&=
x+1,
\\
g_4(x)
&=
1,
&
p_4(x)
&=
1.
&&
\end{align*}
Es folgt die quadratfreie Faktorisierung
$x^5-x^4-5x^3+x^2+8x+4
=
(x-2)^2(x+1)^3
$.
\qedhere
\end{teilaufgaben}
\end{loesung}
